\documentclass{amsart}
% For dealing with references we use the comment environment
\usepackage{verbatim}
\newenvironment{reference}{\comment}{\endcomment}
%\newenvironment{reference}{}{}
\newenvironment{slogan}{\comment}{\endcomment}
\newenvironment{history}{\comment}{\endcomment}

% For commutative diagrams we use Xy-pic
\usepackage[all]{xy}

% We use 2cell for 2-commutative diagrams.
\xyoption{2cell}
\UseAllTwocells

% We use multicol for the list of chapters between chapters
\usepackage{multicol}

% This is generally recommended for better output
\usepackage{lmodern}
\usepackage[T1]{fontenc}

% For cross-file-references

% Package for hypertext links:
\usepackage{hyperref}

% For any local file, say "hello.tex" you want to link to please

% Theorem environments.
%
\theoremstyle{plain}
\newtheorem{theorem}[subsection]{Theorem}
\newtheorem{proposition}[subsection]{Proposition}
\newtheorem{lemma}[subsection]{Lemma}

\theoremstyle{definition}
\newtheorem{definition}[subsection]{Definition}
\newtheorem{example}[subsection]{Example}
\newtheorem{exercise}[subsection]{Exercise}
\newtheorem{situation}[subsection]{Situation}

\theoremstyle{remark}
\newtheorem{remark}[subsection]{Remark}
\newtheorem{remarks}[subsection]{Remarks}

\numberwithin{equation}{subsection}

% Macros
%
\def\lim{\mathop{\mathrm{lim}}\nolimits}
\def\colim{\mathop{\mathrm{colim}}\nolimits}
\def\Spec{\mathop{\mathrm{Spec}}}
\def\Hom{\mathop{\mathrm{Hom}}\nolimits}
\def\Ext{\mathop{\mathrm{Ext}}\nolimits}
\def\SheafHom{\mathop{\mathcal{H}\!\mathit{om}}\nolimits}
\def\SheafExt{\mathop{\mathcal{E}\!\mathit{xt}}\nolimits}
\def\Sch{\mathit{Sch}}
\def\Mor{\mathop{\mathrm{Mor}}\nolimits}
\def\Ob{\mathop{\mathrm{Ob}}\nolimits}
\def\Sh{\mathop{\mathit{Sh}}\nolimits}
\def\NL{\mathop{N\!L}\nolimits}
\def\CH{\mathop{\mathrm{CH}}\nolimits}
\def\proetale{{pro\text{-}\acute{e}tale}}
\def\etale{{\acute{e}tale}}
\def\QCoh{\mathit{QCoh}}
\def\Ker{\mathop{\mathrm{Ker}}}
\def\Im{\mathop{\mathrm{Im}}}
\def\Coker{\mathop{\mathrm{Coker}}}
\def\Coim{\mathop{\mathrm{Coim}}}

% Boxtimes
%
\DeclareMathSymbol{\boxtimes}{\mathbin}{AMSa}{"02}

%
% Macros for moduli stacks/spaces
%
\def\QCohstack{\mathcal{QC}\!\mathit{oh}}
\def\Cohstack{\mathcal{C}\!\mathit{oh}}
\def\Spacesstack{\mathcal{S}\!\mathit{paces}}
\def\Quotfunctor{\mathrm{Quot}}
\def\Hilbfunctor{\mathrm{Hilb}}
\def\Curvesstack{\mathcal{C}\!\mathit{urves}}
\def\Polarizedstack{\mathcal{P}\!\mathit{olarized}}
\def\Complexesstack{\mathcal{C}\!\mathit{omplexes}}
% \Pic is the operator that assigns to X its picard group, usage \Pic(X)
% \Picardstack_{X/B} denotes the Picard stack of X over B
% \Picardfunctor_{X/B} denotes the Picard functor of X over B
\def\Pic{\mathop{\mathrm{Pic}}\nolimits}
\def\Picardstack{\mathcal{P}\!\mathit{ic}}
\def\Picardfunctor{\mathrm{Pic}}
\def\Deformationcategory{\mathcal{D}\!\mathit{ef}}

\setlength{\emergencystretch}{3em}
\begin{document}
\title[Stacks Project, Tag 0ASP]{427 --- Finitely presented modules over generalized valuation rings}
\maketitle
\noindent Copyright (C) 2005--2025 Johan de Jong. Stacks Project authors. Source: \href{https://stacks.math.columbia.edu/tag/0ASP}{Tag 0ASP}.

\noindent Editorial difficulty note (not author text). Difficulty H2: The proof finds a generator whose annihilator remains minimal for every lift of its residue class. It proves the cyclic submodule is relatively divisible by comparing totally ordered principal ideals, splits off that summand and inducts to obtain a cyclic decomposition..

\noindent Editorial provenance note (not author text). History: excerpt from revision \texttt{a0444\allowbreak 6e57e\allowbreak c1fbc\allowbreak 252a8\allowbreak 71afc\allowbreak ec775\allowbreak 2fb28\allowbreak 07b14}, more-algebra.tex, lines 37674--37734. Reference presentation was adapted; original-author context and core support blocks were retained or added. The mathematical wording is unchanged. Licensed under GNU FDL 1.2 or later; no invariant sections or cover texts. See ../licenses/COPYING. Context blocks reproduce author definitions and supporting results; they are not additional counted problems.

\begin{lemma}[Generalized valuation rings]
\label{lemma-generalized-valuation-ring}
\begin{reference}
\href{https://stacks.math.columbia.edu/bibliography}{Bibliography: Warfield-Decomposition}
\end{reference}
Let $R$ be a nonzero ring. The following are equivalent
\begin{enumerate}
\item For $a, b \in R$ either $a$ divides $b$ or $b$ divides $a$.
\item Every finitely generated ideal is principal and $R$ is local.
\item The set of ideals of $R$ is linearly ordered by inclusion.
\end{enumerate}
This holds in particular if $R$ is a valuation ring.
\end{lemma}

\begin{lemma}
\label{lemma-generalized-valuation-ring-modules}
\begin{reference}
\href{https://stacks.math.columbia.edu/bibliography}{Bibliography: Warfield-Decomposition (Theorem 1)}
\end{reference}
Let $R$ be a ring satisfying the equivalent conditions of
Lemma \ref{lemma-generalized-valuation-ring}.
Then every finitely presented $R$-module
is isomorphic to a finite direct sum of modules of the form $R/fR$.
\end{lemma}

\begin{proof}
Let $M$ be a finitely presented $R$-module. We will use all
the equivalent properties of $R$ from
Lemma \ref{lemma-generalized-valuation-ring}
without further mention. Denote $\mathfrak m \subset R$
the maximal ideal and $\kappa = R/\mathfrak m$ the residue field.
Let $I \subset R$ be the annihilator of $M$.
Choose a basis $y_1, \ldots, y_n$ of the finite dimensional
$\kappa$-vector space $M/\mathfrak m M$. We will argue
by induction on $n$.

\medskip\noindent
By Nakayama's lemma any collection of elements $x_1, \ldots, x_n \in M$
lifting the elements $y_1, \ldots, y_n$ in $M/\mathfrak m M$
generate $M$, see Algebra, Lemma \href{https://stacks.math.columbia.edu/tag/00DV}{lemma NAK (Tag 00DV)}.
This immediately proves the base case $n = 0$ of the induction.

\medskip\noindent
We claim there exists an index $i$ such that for any choice
of $x_i \in M$ mapping to $y_i$ the annihilator of $x_i$ is $I$.
Namely, if not, then we can choose $x_1, \ldots, x_n$ such
that $I_i = \text{Ann}(x_i) \not = I$ for all $i$. But as
$I \subset I_i$ for all $i$, ideals being totally ordered
implies $I_i$ is strictly bigger than $I$ for $i = 1, \ldots, n$,
and by total ordering once more we would see that
$\text{Ann}(M) = I_1 \cap \ldots \cap I_n$ is bigger than $I$
which is a contradiction. After renumbering we may assume that
$y_1$ has the property: for any $x_1 \in M$ lifting $y_1$
the annihilator of $x_1$ is $I$.

\medskip\noindent
We set $A = Rx_1 \subset M$. Consider the exact sequence
$0 \to A \to M \to M/A \to 0$. Since $A$ is finite, we see that
$M/A$ is a finitely presented $R$-module
(Algebra, Lemma \href{https://stacks.math.columbia.edu/tag/0519}{lemma extension (Tag 0519)}) with fewer generators.
Hence $M/A \cong \bigoplus_{j = 1, \ldots, m} R/f_jR$ by induction.
On the other hand, we claim that $A \to M$ satisfies the property:
if $f \in R$, then $fA = A \cap fM$. The inclusion $fA \subset A \cap fM$
is trivial. Conversely, if $x \in A \cap fM$, then $x = gx_1 = f y$
for some $g \in R$ and $y \in M$. If $f$ divides $g$, then $x \in fA$
as desired. If not, then we can write $f = hg$ for some $h \in \mathfrak m$.
The element $x'_1 = x_1 - hy$ has annihilator $I$ by the previous
paragraph. Thus $g \in I$ and we see that $x = 0$ as desired.
The claim and Lemma \href{https://stacks.math.columbia.edu/tag/0ASM}{lemma characterize PD modules (Tag 0ASM)}
imply the sequence $0 \to A \to M \to M/A \to 0$ is split
and we find $M \cong A \oplus \bigoplus_{j = 1, \ldots, m} R/f_jR$.
Then $A = R/I$ is finitely presented (as a summand of $M$)
and hence $I$ is finitely generated, hence principal.
This finishes the proof.
\end{proof}
\end{document}
